\scanpage{16}
avec des $A_i^E(\unsure{\zeta|\xi})$ $(0\leq i\leq n-1)$ \unsure{uniquemt|uniquement} déterminés.
On trouve ainsi des hom
\[
 A_i^E\in\Gamma(\mathrm V(E)^0,\deleted{\underline O_{\mathrm V(E)^0})}
 \qquad\text{i.e.}\qquad A_i^E:\mathrm V(E)^0\longrightarrow G_{aS}
\]
uniquemt déterminés par la condition précédente.
On peut les interpréter comme des éléments de
\[
 \underline{\mathrm{Sym}}^*(E)_{\Delta_n^E}
\]
\deleted{$A_i$ est homogène de degré $p$}

Ceci dit, on trouve que les $A_i$ sont des sections de $\underline{\mathrm{Sym}}(E)$, i.e. sont définis sur
$\mathrm V(E)$ tt entier, bien entendu ils y sont caractérisés par l'identité précédente.
\deleted{\unsure{$p^{n}$} \unsure{est}}

On a \unsure{encore|même}
\[
 A_0^E=(\Delta_n^E)^{p-1}\deleted{(\unsure{-1})^n}(\pm1)
\]
[celà a un sens, grâce à :
\[
 \begin{aligned}
 &(\textstyle\bigwedge^n\check E)^{p-1}\simeq\underline O_S\\
 \text{i.e.}\quad&(\textstyle\bigwedge^n\check E)^p=\textstyle\bigwedge^n\check E^{(p)}\simeq
 (\textstyle\bigwedge^n\check E)\quad\big]
 \end{aligned}
\]

\uline{NB} Supposons que $\zeta$ est une section de $\mathrm V(E)^0$, \deleted{Alors}
donc les $\zeta^{p^i}$ $(0\leq i\leq n-1)$ font une base de $\check E$, qui s'identifie donc à
$\underline O_S^{n}$, avec la \deleted{\unsure{loi|forme|choix}} de \unsure{puiss.} \unsure{$p$ième}
$e_0^{(p)}=e_1,e_1^{(p)}=e_2,\ldots,e_{n-2}^{(p)}=e_{n-1}$, et
% archive sheet number: 15
